\subsection{Exploitability}

In einem Zweipersonen-Nullsummenspiel ist die exploitability einer Strategie $\sigma_i$ definiert als
\begin{equation}
\varepsilon_i(\sigma_i) = v_i - u_i(\sigma_i, b_{-i}(\sigma_i)),
\end{equation}
wobei $v_i$ der Spielwert für Spieler $i$ ist und $b_{-i}(\sigma_i)$ die best response des Gegners gegen $\sigma_i$ darstellt.
Eine Strategie ist optimal, wenn ihre exploitability null ist.

Für Strategien, die in beiden Positionen gespielt werden, wird die exploitability als Durchschnitt der best response Werte aus beiden Positionen berechnet:
\begin{equation}
\varepsilon(\sigma) = \frac{u_2(\sigma_1, b_2(\sigma_1)) + u_1(b_1(\sigma_2), \sigma_2)}{2}.
\end{equation}

Diese Metrik ermöglicht es, die Worst-Case-Performance von Strategien in großen Spielen zu bewerten, wo eine vollständige Traversierung des game trees nicht möglich ist.
